\documentclass[10pt,a4paper]{amsart}
\usepackage[margin=19mm]{geometry}
\usepackage{amsmath,amssymb,mathtools}
\usepackage[scr=rsfs,cal=euler]{mathalfa}
\usepackage{xfrac}
\usepackage[unicode,hidelinks,bookmarks=false]{hyperref}
\newcommand{\ie}{i.e.\ }
\newcommand{\cf}{cf.\ }
\newcommand{\resp}{respectively\ }
\newcommand{\Subsection}{\subsection}
\providecommand{\nobreakdash}{\nobreak\hbox{-}}
\setlength{\emergencystretch}{2em}
\multlinegap0pt
\usepackage{bm}
\usepackage{bbm}
\usepackage{dsfont}
\usepackage{xspace}
\usepackage{cancel}
\usepackage{mathtools}
\usepackage{amsthm}
\makeatletter
\@ifundefined{thm}{\newtheorem{thm}{Theorem}}{}
\@ifundefined{defi}{\newtheorem{defi}[thm]{Definition}}{}
\makeatother
\newcommand {\emptycomment}[1]{}
\newcommand{\Der}{\mathrm{Der}}
\newcommand{\Hom}{\mathrm{Hom}}
\newcommand{\LD}{{\rm LieDer }}
\newcommand{\ad}{\mathrm{ad}}
\newcommand{\dM}{\mathrm{d}}
\newcommand{\frkk}{\mathfrak k}
\newcommand{\frkz}{\mathfrak z}
\newcommand{\g}{\mathfrak g}
\newcommand{\h}{\mathfrak h}
\newcommand{\lon }{\,\rightarrow\,}
\title[Non-abelian Extensions of Lie algebras with derivations]{Non-abelian Extensions of Lie algebras with derivations}
\author{Jun Jiang, Kanghe Xu}
\date{Source: arXiv:2604.26276v3 (CC BY 4.0)}
\begin{document}
\maketitle

\noindent\emph{Source and licence.} Non-abelian Extensions of Lie algebras with derivations. Version arXiv:2604.26276v3 (CC BY 4.0), distributed under the Creative Commons Attribution 4.0 International licence, as recorded in the arXiv licence metadata for this exact version and shown on the abstract page https://arxiv.org/abs/2604.26276v3. Full rights notice: licenses/arxiv-2604.26276-CC-BY-4.0.txt.\par\smallskip

\noindent\textit{Editorial scope.} Counted unit: the Thm whose source label is \texttt{gdk} in the article (source0.tex lines 1120--1125, source label \texttt{gdk}), delivered together with the article's own complete original proof of it (source0.tex lines 1126--1204, one \texttt{proof} environment of 79 lines). No mathematics is written, repaired or re-derived: statement and proof are the article's own text line for line. Required context retained uncounted: nc1 (lines 390--399); nc2 (lines 392--397); nc3 (lines 392--397); nc4 (lines 392--397); corres (lines 539--542). Every definition, display and label of those lines is retained unchanged. Omitted from this package and not redistributed: 1--389, 398--538, 543--1119, 1205--1836. Nothing inside a retained excerpt is omitted or altered: every retained body is delivered exactly as the article's lines are written, so no omission receipt of any kind accompanies this package. Citations: the retained excerpts carry no citation command at all, so no bibliography entry of the article is transcribed and no reference is redistributed. Reference adaptation: none was needed and none was made; every reference inside the delivered unit resolves inside it, so no unresolved reference or citation is produced. Printed slips kept literal: no printed word, symbol, hypothesis or number of the counted statement or of its proof was altered. Layout and class: the article's own class is \texttt{amsart} with packages including cancel, amsmath, amsfonts, amssymb, xy, mathrsfs, bbding, txfonts, amscd, enumitem, ifpdf, hyperref, tikz, srcltx; the wrapper is the lane's pinned \texttt{amsart} editorial wrapper at 10pt on A4, which restates the theorem-like environments the retained text needs and adds the article's own macro definitions that the retained text needs (\texttt{\textbackslash Der}, \texttt{\textbackslash Hom}, \texttt{\textbackslash LD}, \texttt{\textbackslash ad}, \texttt{\textbackslash dM}, \texttt{\textbackslash emptycomment}, \texttt{\textbackslash frkk}, \texttt{\textbackslash frkz}, \texttt{\textbackslash g}, \texttt{\textbackslash h}, \texttt{\textbackslash lon}), with extra wrapper packages (bm, bbm, dsfont, xspace, cancel, mathtools). The delivered numbering is the unit's own. Assets: no retained span contains a figure environment, a graphics-inclusion command, a PDF-page inclusion, a table or a TikZ picture; no image file is copied into this package, the package declares no asset dependency and no image file is redistributed with it. Licence: version arXiv:2604.26276v3 of this article is distributed under the Creative Commons Attribution 4.0 International licence, as recorded in the arXiv licence metadata for this exact version and shown on the abstract page https://arxiv.org/abs/2604.26276v3. A full rights notice is delivered at licenses/arxiv-2604.26276-CC-BY-4.0.txt. Characters: every retained excerpt is pure ASCII, so the delivered engine, XeLaTeX with its default Latin Modern text font, reports no missing character.
\begin{defi}
Let $(\g, D)$ and $(\h, K)$ be two {\rm LieDer} pairs. A {\bf non-abelian $2$-cocycle} of $(\g, D)$ with values in $(\h, K)$ is a triple $(\varrho, \omega, \chi)$ of linear maps $\varrho: \g\lon\Der(\h),~\omega:\wedge^2\g\lon\h$ and $\chi:\g\lon\h$ such that for all $x, y, z\in\g$ and $u, v\in\h$ the following equations hold:
\begin{eqnarray}
\label{nc1}\varrho([x, y]_\g)u&=&\varrho(x)\varrho(y)u-\varrho(y)\varrho(x)u-[\omega(x, y), u]_\h,\\
\label{nc2}\varrho(x)\omega(y, z)+\varrho(y)\omega(z, x)&+&\varrho(z)\omega(x, y)-\omega([x, y]_\g, z)-\omega([y, z]_\g, x)-\omega([z, x]_\g, y)=0,\\
\label{nc3}K(\varrho(x)v)&=&\varrho(D(x))v+\varrho(x)K(v)+[\chi(x), v]_\h,\\
\label{nc4}K(\omega(x, y))+\chi([x, y]_\g)&=&\varrho(x)\chi(y)-\varrho(y)\chi(x)+\omega(D(x), y)+\omega(x, D(y)).
\end{eqnarray}
Denote by $Z^2_{nab}(\g, D;\h, K)$ the set of non-abelian $2$-cocycles.
\end{defi}

\begin{thm}\label{corres}
Let $(\g, D)$ and $(\h, K)$ be two {\rm LieDer} pairs. Then the non-abelian extensions of $(\g, D)$ by $(\h, K)$ are classified by the second non-abelian cohomology $H_{nab}^2
(\g, D; \h, K)$.
\end{thm}

\begin{thm}\label{gdk}
With the above notations, the set $\mathrm{Ext}(\g,D;\h,K)_{\frkk}$ is isomorphic to $H^2_{\LD}(\g,D;\frkz(\h),K)_{\frkk}$. Moreover, we have 
    \begin{equation*}
      {\rm{\bf Ext_{nab}}}(\g,D;\h,K)=\bigsqcup_{\text{Integrable}~~(\g,D)\text{-kernels}}H^2_{\LD}(\g,D; \frkz(\h), K)_{\frkk}.
    \end{equation*}
    \end{thm}

    \begin{proof}
Fix an element $[(\hat{\g}, \hat{D})]$ in $\rm{Ext}(\g, D;\h,K)_\frkk$, it means that $[(\hat{\g}, \hat{D})]$ is the isomorphism class of $(\hat{\g}, \hat{D})$, where $(\hat{\g}, \hat{D})$ is a non-abelian extension of $(\g,D)$ by $(\h,K)$
whose associated $(\g, D)$-kernel is $\frkk$. 
By Theorem \ref{corres}, the non-abelian extensions of $(\g, D)$ by $(\h, K)$ are classified by the second non-abelian cohomology $H^2_{nab}(\g,D; \h,K)$. Let $[(\varrho, \omega, \chi)]$ be the element in $H^{2}_{nab}(\g, D;\h,K)$ corresponding to $[(\hat{\g}, \hat{D})]$. Now, let $[(\tilde{\g}, \tilde{D})]\in{\rm Ext}(\g, D;\h,K)_\frkk$ and $[(\tilde{\g}, \tilde{D})]\neq [(\hat{\g}, \hat{D})]$. Denote by $[(\tilde{\varrho}, \tilde{\omega}, \tilde{\chi})]$ the element in $H^{2}_{nab}(\g,D; \h,K)$ corresponding to $[(\tilde{\g}, \tilde{D})]$. Then $\overline{\varrho}=\overline{\tilde{\varrho}}=\frkk$, which means that there exists a linear map $\tau:\g\lon\h$ such that
$$\varrho(x)-\tilde{\varrho}(x)=\ad\tau(x), \quad \forall x\in\g.$$
Define linear maps $\omega^*:\wedge^2\g\lon\h$ and $\chi^*:\g\lon\h$ by
\begin{eqnarray}
\label{H1}\omega^*(x, y)&=&\tilde{\omega}(x, y)+\tilde{\varrho}(x)\tau(y)-\tilde{\varrho}(y)\tau(x)+[\tau(x), \tau(y)]_\h-\tau([x, y]_\g),\\
\label{H2}\chi^*(x)&=&\tilde{\chi}(x)+K(\tau(x))-\tau(D(x)),
\end{eqnarray}
for all $x, y\in\g$. Then $[(\tilde{\varrho}, \tilde{\omega}, \tilde{\chi})]=[(\varrho, \omega^*, \chi^*)]$. By \eqref{nc1} and \eqref{nc3}, it implies that
$$
[\omega(x, y)-\omega^*(x, y), u]_\h=0,\quad [\chi(x)-\chi^*(x), u]_\h=0,
$$
for all $x, y\in\g, u\in\h$,
which means that the map $\eta=\omega-\omega^*\in\Hom(\wedge^2\g, \frkz(\h))$ and $\theta=\chi-\chi^*\in\Hom(\g, \frkz(\h))$. By \eqref{nc2} and \eqref{nc4}, we have
$$
\dM^{CE}_{\rho_\frkk}(\eta)=0, \quad \dM^{CE}_{\rho_\frkk}(\theta)+\delta(\eta)=0,
$$
thus, $\partial_{\rho_\frkk}(\eta, \theta)=0$, i.e.,$(\eta, \theta)$ is a $2$-cocycle of the complex $(C^{*}_{\mathrm{LieDer}}(\g, \frkz(\h)), \partial_{\rho_\frkk})$. 

Suppose that $[(\varrho', \omega', \chi')]=[(\varrho, \omega, \chi)]$ and $
[(\tilde{\varrho}', \tilde{\omega}', \tilde{\chi}')]=[(\tilde{\varrho}, \tilde{\omega}, \tilde{\chi})]$ such that $\overline{\varrho}=\overline{\varrho'}=\overline{\tilde{\varrho}}=\overline{\tilde{\varrho}'}=\frkk$. Then there exist linear maps $\tau', \tilde{\tau}':\g\lon\h$ such that $(\varrho', \omega', \chi')$ is equivalent to $(\varrho, \omega, \chi)$, and $(\tilde{\varrho}', \tilde{\omega}', \tilde{\chi}')$ is equivalent to $(\tilde{\varrho}, \tilde{\omega}, \tilde{\chi})$.
Moreover, there exists $\tau^{*}{'}:\g\lon\h$ such that
\emptycomment{
\begin{eqnarray*}
\left\{\begin{array}{rcl}
\varrho(x)u-\varrho'(x)u&=&[\tau'(x), u]_\h,\\
\omega(x, y)-\omega'(x, y)&=&\varrho'(x)\tau'(y)-\varrho'(y)\tau'(x)+[\tau'(x), \tau'(y)]_\h-\tau'([x, y]_\g),\\
\chi(x)-\chi'(x)&=&K(\tau'(x))-\tau'(D(x)),
\end{array}\right.
\end{eqnarray*}
\begin{eqnarray*}
\left\{\begin{array}{rcl}
\tilde{\varrho}(x)u-\tilde{\varrho}'(x)u&=&[\tilde{\tau}'(x), u]_\h,\\
\tilde{\omega}(x, y)-\tilde{\omega}'(x, y)&=&\tilde{\varrho}'(x)\tilde{\tau}'(y)-\tilde{\varrho}'(y)\tilde{\tau}'(x)+[\tilde{\tau}'(x), \tilde{\tau}'(y)]_\h-\tilde{\tau}'([x, y]_\g),\\
\tilde{\chi}(x)-\tilde{\chi}'(x)&=&K(\tilde{\tau}'(x))-\tilde{\tau}'(D(x)),
\end{array}\right.
\end{eqnarray*}
for all $x, y\in\g, u\in\h$, and}
$$\varrho'(x)-\tilde{\varrho}'(x)=\ad\tau^{*}{'}(x), \quad\forall x\in\g.$$ Then $[(\varrho', \omega^{*}{'}, \chi^{*}{'})]=[(\tilde{\varrho}', \tilde{\omega}', \tilde{\chi}')]$, where
\begin{eqnarray}
\label{H3}\omega^{*}{'}(x, y)&=&\tilde{\omega}'(x, y)+\tilde{\varrho}'(x)\tau^{*}{'}(y)-\tilde{\varrho}'(y)\tau^{*}{'}(x)+[\tau^{*}{'}(x), \tau^{*}{'}(y)]_\h-\tau^{*}{'}([x, y]_\g),\\
\label{H4}\chi^{*}{'}(x)&=&\tilde{\chi}'(x)+K(\tau^{*}{'}(x))-\tau^{*}{'}(D(x)).
\end{eqnarray}
Moreover, the maps $\eta'=\omega'-\omega^{*}{'}$ and $\theta'=\chi'-\chi^{*}{'}$ satisfy $\partial_{\rho_\frkk}(\eta', \theta')=0$. By \eqref{H1}, \eqref{H2}, \eqref{H3} and \eqref{H4}, it implies that
\begin{eqnarray*}
(\eta-\eta')(x, y)=\omega(x, y)-\omega'(x, y)-\omega^{*}(x, y)+\omega^{*}{'}(x, y) =\dM^{CE}_{\rho_\frkk}(\tau'+\tau^{*}{'}-\tilde{\tau}{'}-\tau)(x, y),
\end{eqnarray*}
and 
\begin{eqnarray*}
(\theta-\theta')(x)=\chi(x)-\chi'(x)-\chi^*(x)+\chi^{*}{'}(x)=-\delta(\tau'+\tau^{*}{'}-\tilde{\tau}'-\tau)(x),
\end{eqnarray*}
which means that $(\eta, \theta)-(\eta', \theta')=\partial_{\rho_{\frkk}}(\tau'+\tau^{*}{'}-\tilde{\tau}'-\tau)$, i.e.,$[(\eta, \theta)]=[(\eta', \theta')]\in H^{2}(\g, \frkz(\h))_{\frkk}$. Therefore the map  
$\Xi:{\rm Ext}(\g,D; \h, K)_\frkk\lon H_{\LD}^{2}(\g,D; \frkz(\h),K)_\frkk$ defined by
$$
\Xi([(\tilde{\g}, \tilde{D})])=[(\eta, \theta)], 
$$
is well defined, and $\Xi([(\hat{\g}, \hat{D})])=0$.

Assume that $\Xi([(\tilde{\g}, \tilde{D})])=0,$ there exists a linear map $l:\g\lon\frkz(\h)$ such that $\partial_{\rho_{\frkk}}(l)=(\eta, \theta)$. By direct calculation, 
\emptycomment{
	Then
\begin{eqnarray*}
\left\{\begin{array}{rcl}
\varrho(x)u-\tilde{\varrho}(x)u&=&[\tau(x)+l(x), u]_\h,\\
\omega(x, y)-\tilde{\omega}(x, y)&=&\tilde{\varrho}(x)(\tau+l)(y)-\tilde{\varrho}(y)(\tau+l)(x)+[(\tau+l)(x), (\tau+l)(y)]_\h-(l+\tau)([x, y]_\g),\\
\chi(x)-\tilde{\chi}(x)&=&K((\tau+l)(x))-(\tau+l)(D(x)),
\end{array}\right.
\end{eqnarray*}
which means that }$[(\varrho, \omega, \chi)]=[(\tilde{\varrho}, \tilde{\omega}, \tilde{\chi})]$, i.e., $[(\tilde{\g}, \tilde{D})]=[(\hat{\g}, \hat{D})]$. Thus $\Xi$ is injective. 

For any $[(\eta, \theta)]\in H^2_{\LD}(\g,D; \frkz(\h),K)_\frkk$, i.e., $\partial_{\rho_\frkk}(\eta, \theta)=0$. By direct calculation, $(\varrho, \omega-\eta, \chi-\theta)$ is a non-abelian $2$-cocycle. Denote by $(\bar{\g}, \bar{D})$ the non-abelian extension of $(\g, D)$ by $(\h, K)$ corresponding to $(\varrho, \omega-\eta, \chi-\theta)$. We have $\Xi([(\bar{\g}, \bar{D})])=[(\eta, \theta)]$, which means that $\Xi$ is surjective. Thus,
the set $\mathrm{Ext}(\g,D;\h,K)_{\frkk}$ is isomorphic to $H^2_{\LD}(\g,D;\frkz(\g),K)_{\frkk}$, which implies 
\begin{equation*}
      {\rm{\bf Ext_{nab}}}(\g,D;\h,K)=\bigsqcup_{\text{Integrable}(\g,D)\text{-kernels}}H^2_{\LD}(\g,D;\frkz(\h),K)_{\frkk}.
    \end{equation*}
We finish the proof.
    \end{proof}

\end{document}
